\documentclass[11pt]{article}
\usepackage[a4paper,margin=1in]{geometry}
\usepackage{amsmath}
\usepackage{booktabs}
\usepackage[T1]{fontenc}

\title{Evidence-Aware Research Communication with Retrieval}
\author{SARAL Prototype Demonstration}
\date{September 2026}

\begin{document}
\maketitle

\begin{abstract}
Researchers often need to explain a paper to audiences with different levels of technical knowledge. This short demonstration paper describes a retrieval-backed communication workflow. The workflow maps generated statements to source chunks so that readers can inspect the evidence behind each statement.
\end{abstract}

\section{Problem}
Summaries are useful only when a reader can tell which evidence supports them. A policymaker may need a concise explanation of public value, while graduate students may need methods, assumptions, and equations. A single fixed summary does not serve both audiences well.

\section{Method}
The system splits a paper into page-aware chunks and ranks the chunks against an audience-specific request. A generator uses only the selected evidence to create slide bullets, speaker notes, and a script. Each generated item stores one or more chunk identifiers and page numbers.

We measure citation coverage as:
\[
\mathrm{CitationCoverage} =
\frac{\text{number of generated claims with at least one citation}}
{\text{total number of generated claims}}.
\]

We also use lexical grounding as a lightweight factuality proxy. For a claim $g$ and its retrieved evidence $e$, the score is the fraction of unique claim tokens that also occur in the evidence:
\[
\mathrm{Grounding}(g,e) = \frac{|\mathrm{tokens}(g) \cap \mathrm{tokens}(e)|}{|\mathrm{tokens}(g)|}.
\]

\section{Illustrative Findings}
In a small internal demonstration, evidence-linked outputs made it easier to verify statements during review. Graduate-student outputs retained the evaluation definitions, while policymaker outputs foregrounded transparency and accountability. The workflow should not treat overlap metrics as proof of factuality; human review remains necessary for high-stakes use.

\begin{table}[h]
\centering
\begin{tabular}{lll}
\toprule
Audience & Primary need & Recommended output \\
\midrule
Policymakers & Public value and traceability & 90-second plain-English script \\
Graduate students & Methods and equations & Technical five-slide talk \\
Press & Clear, quotable framing & 30-second press summary \\
\bottomrule
\end{tabular}
\caption{Example audience-specific communication requests.}
\end{table}

\section{Limitations}
Sparse retrieval can miss relevant paraphrases, and a generated citation does not itself guarantee that the claim is correct. A production system should add semantic retrieval, independent citation checking, safety review, feedback collection, and access controls.

\section{Conclusion}
Retrieval-backed generation can make research communication more adaptable while preserving a clear path back to the underlying paper. The strongest demonstration is an iterative edit that changes the wording for an audience without losing the evidence links.

\end{document}

